\chapter{Work and Mechanical Energy}\label{ch:g12:work-and-energy}

Tip a marble into a salad bowl: it dives, overshoots the bottom,
climbs to its starting height, hesitates, turns back --- and every
reversal can be predicted from one curve (the bowl's profile) and one
horizontal line (the marble's energy), no equation of motion solved.
Last year priced forces on straight lines (\cref{ch:g11:work-of-force})
and balanced kinetic against potential energy
(\cref{ch:g11:mechanical-energy}); this year the bookkeeping goes
everywhere: curved paths, springs, landscapes at a glance.

\section{Work along any path}

\begin{definition}[Work along a path]\label{def:g12:work-and-energy:path}
Let a force $\vect F$ act on a body moving along any path from $A$ to
$B$: cut the path into displacements $\vect{\dd\ell}$ so short that
each is straight and $\vect F$ barely changes along it. Each step earns
the small work $\scal{F}{\dd\ell} = F\,\dd\ell\cos\theta$
(\cref{ch:g11:work-of-force}), and the
\emph{work along the path}\index{work (of a force)!along a path} is
their sum, $W_{AB}(\vect F) = \sum \scal{F}{\dd\ell}$ --- which on a
straight path with constant force collapses to last year's
$F \times AB \times \cos\theta$.
\end{definition}

\begin{omfigure}
\begin{tikzpicture}[font=\small]
  \draw[omExo!60, thick] plot [smooth, tension=0.8] coordinates {(0,0.3) (1.2,1.15) (2.6,1.3) (4.0,0.7) (5.3,1.05)};
  \draw[omMeth, thick] (0,0.3) -- (1.2,1.15) -- (2.6,1.3) -- (4.0,0.7) -- (5.3,1.05);
  \foreach \p in {(0,0.3),(1.2,1.15),(2.6,1.3),(4.0,0.7),(5.3,1.05)} \fill \p circle (1.3pt);
  \node[below] at (0,0.3) {$A$}; \node[below right] at (5.3,1.05) {$B$};
  \draw[-Stealth, omThm, very thick] (2.6,1.3) -- (3.44,0.94) node[below left=-2pt] {$\vect{\dd\ell}$};
  \draw[-Stealth, omDef, very thick] (2.6,1.3) -- ++(22:1.5) node[above] {$\vect F$};
  \draw (2.6,1.3) ++(-23:0.75) arc[start angle=-23, end angle=22, radius=0.75];
  \node at ($(2.6,1.3)+(0:1.0)$) {$\theta$};
\end{tikzpicture}

{\small A curved path is a broken line of many short steps: each
$\vect{\dd\ell}$ earns $F\,\dd\ell\cos\theta$; the work is the sum.}
\end{omfigure}

\begin{proposition}[Work of the weight, any path]\label{prop:g12:work-and-energy:weight}
Along \emph{any} path from $A$ (altitude $z_A$) to $B$ (altitude
$z_B$), the weight of a mass $m$ does the work
$W_{AB}(\vect P) = mg\,(z_A - z_B)$: the drop enters, never the route.
\end{proposition}

\begin{proof}
Each short step contributes $mg \times (\text{small drop})$ (last
year's straight-segment computation); the drops telescope to
$z_A - z_B$. The curved case, admitted then, is delivered.
\end{proof}

\begin{definition}[Conservative and non-conservative forces]\label{def:g12:work-and-energy:conservative}
A force is \emph{conservative}\index{conservative force} if its work
between two points is the same along every path (zero on every round
trip), \emph{non-conservative}\index{non-conservative force}
otherwise. Weight and spring force (next section): conservative.
Sliding friction is not: opposite the motion, it does $W = -fL$ on a
path of length $L$ --- a toll per metre, never refunded.
\end{definition}

\begin{example}[Two trails to the hut]\label{ex:g12:work-and-energy:trails}
A \qty{75}{kg} hiker gains \qty{300}{m} by a steep \qty{800}{m} trail
or \qty{2.4}{km} of switchbacks. The weight charges
$-75 \times 9.81 \times 300 \approx \qty{-2.2e5}{J}$ on both; a
\qty{30}{N} friction force, \qty{-2.4e4}{J} on one, \qty{-7.2e4}{J}
on the other. The weight refunds on the descent; friction's bill is
gone as heat, both ways.
\end{example}

\section{Potential energy}

\begin{definition}[Potential energy]\label{def:g12:work-and-energy:ep}
To every conservative force one attaches a
\emph{potential energy}\index{potential energy!from conservative work}
$E_p$, a number depending on position alone, defined by
$W_{AB} = E_p(A) - E_p(B)$: work delivered is potential energy spent.
Fixing $E_p = 0$ at a chosen reference point fixes it everywhere; the
choice shifts all values by a constant that cancels from every
difference. For the weight this gives back last year's $E_p = mgz$,
now valid along any path. Friction can own no $E_p$: a route-dependent
bill is no difference of endpoint numbers --- only conservative forces
store retrievably, exactly what the name conserves.
\end{definition}

\begin{proposition}[Elastic potential energy]\label{prop:g12:work-and-energy:elastic}
A spring of stiffness $k$ (\unit{N/m}) stretched or compressed by $x$
pulls back with the force $F = kx$
(\cref{ch:g12:oscillators-and-time}); the work it delivers returning
to rest --- its
\emph{elastic potential energy}\index{potential energy!elastic} --- is
$E_p = \tfrac12\, k x^2$.
\end{proposition}

\begin{proof}
The restoring force acts along the motion, so $\sum F\,\dd\ell$ is
the \emph{area under the $F$--$x$ graph}: a triangle of base $x$ and
height $kx$, area $\tfrac12 x \times kx$. The force at each stretch
is the same out and back: work fixed by the endpoints, spring force
conservative, $E_p$ well defined.
\end{proof}

\begin{omfigure}
\begin{tikzpicture}
\begin{axis}[omaxis, width=8.4cm, height=4.8cm, xmin=0, xmax=0.095,
    ymin=0, ymax=21,
    xlabel={stretch $x$ (\unit{m})}, ylabel={force $F$ (\unit{N})}]
  \addplot[draw=none, fill=omDef!18, domain=0:0.06, samples=2] {200*x} \closedcycle;
  \addplot[omThm, thick, domain=0:0.09, samples=2] {200*x};
  \draw[dashed, omExo] (axis cs:0,12) -- (axis cs:0.06,12) -- (axis cs:0.06,0);
  \node[font=\small, anchor=west] at (axis cs:0.061,6.5) {$kx$};
  \node[font=\small] at (axis cs:0.044,3.2) {$\tfrac12 kx^2$};
  \node[font=\small, anchor=south, omThm] at (axis cs:0.082,17) {$F = kx$};
\end{axis}
\end{tikzpicture}

{\small The spring's force grows with the stretch; the stored energy
is the area under the line: a triangle, $\tfrac12 \times x \times kx$.}
\end{omfigure}

\section{The two energy theorems}

\begin{theorem}[Kinetic-energy theorem]\label{thm:g12:work-and-energy:ketheorem}\index{kinetic-energy theorem}
For any motion of a body of mass $m$ from $A$ to $B$ --- curved,
looped, three-dimensional ---
\[
\Delta E_k = E_k(B) - E_k(A) = \sum W_{AB}(\vect F),
\]
the sum running over all forces acting on the body.
\end{theorem}

\begin{proof}
Since $v^2 = \scal{v}{v}$, the product rule of this year's mathematics
gives the derivative of $v^2$ as $2\,\scal{a}{v}$, so
$\dd E_k/\dd t = m\,\scal{a}{v} = \scal{F_{\text{tot}}}{v}$ by
Newton's second law (\cref{ch:g12:newtons-laws}). In a short $\dd t$
the body moves $\vect{\dd\ell} = \vect v\,\dd t$: $E_k$ gains the
small work of \cref{def:g12:work-and-energy:path}; summing the steps
sums the works --- last year's admitted statement, earned in full.
\end{proof}

\begin{proposition}[Instantaneous power]\label{prop:g12:work-and-energy:power}
At every instant a force $\vect F$ on a body of velocity $\vect v$
delivers the power $P = \scal{F}{v} = Fv\cos\theta$, and
$\dd E_k/\dd t$ is the sum of the powers of all forces: last year's
constant-velocity formula, exact at each instant of any motion.
\end{proposition}

\begin{proof}
Read off the previous proof: each force feeds
$\scal{F}{\dd\ell} = \scal{F}{v}\,\dd t$ into $E_k$ during $\dd t$.
\end{proof}

\begin{theorem}[Mechanical-energy theorem]\label{thm:g12:work-and-energy:emtheorem}\index{mechanical-energy theorem}
Let $E_m = E_k + E_p$, where $E_p$ collects the potential energies of
all the conservative forces at work. Then, along any path,
\[
\Delta E_m = W_{AB}(\text{non-conservative forces}),
\]
and $E_m$ is conserved whenever the non-conservative forces do no work.
\end{theorem}

\begin{proof}
Split the kinetic-energy theorem: $\Delta E_k = W_{\text{cons}} +
W_{\text{nc}} = -\Delta E_p + W_{\text{nc}}$; move $\Delta E_p$ across.
\end{proof}

\begin{example}[The pendulum, honestly this time]\label{ex:g12:work-and-energy:pendulum}
A pendulum bob on a wire of length $L = \qty{2.5}{m}$ is released at
rest at $\qty{45}{\degree}$ from the vertical. The tension,
perpendicular to the motion, is powerless; the weight is conservative:
$E_m$ is conserved \emph{on the curved arc}, which last year could
only be checked. The bob starts $h = L(1 - \cos\qty{45}{\degree}) =
\qty{0.73}{m}$ up, so at the bottom $v = \sqrt{2gh} \approx
\qty{3.8}{m/s}$.
\end{example}

\begin{method}[Energy audit of any motion]\label{met:g12:work-and-energy:audit}
\begin{enumerate}
  \item Choose the system and list the forces on it.
  \item Sort: perpendicular to the motion $\to$ no work; conservative
        $\to$ an $E_p$ in $E_m$; the rest $\to$ $W_{\text{nc}}$
        (sliding friction: $-fL$, $L$ = path length).
  \item Write $\Delta E_m = W_{\text{nc}}$ between the data point and
        the question point; solve. Energy answers ``how fast, how
        high''; only time and acceleration need Newton's laws.
\end{enumerate}
\end{method}

\section{Energy diagrams}

\begin{definition}[Energy diagram, turning points]\label{def:g12:work-and-energy:diagram}
For a motion along one axis $x$ with conserved $E_m$, the
\emph{energy diagram}\index{energy diagram} plots the curve $E_p(x)$
and the horizontal line $E_m$. Since $E_k = E_m - E_p \geq 0$, the
motion is confined to where the curve lies below the line; the gap
between them is $E_k$, read directly. At a
\emph{turning point}\index{turning point}, where curve meets line,
$v = 0$ and the motion reverses.
\end{definition}

\begin{omfigure}
\begin{tikzpicture}
\begin{axis}[omaxis, width=9.0cm, height=5.4cm,
    xlabel={$x$}, ylabel={energy}, xtick=\empty, ytick=\empty,
    xmin=-2.25, xmax=2.25, ymin=0, ymax=2.45]
  \addplot[omDef, thick, domain=-2.05:2.05, samples=80] {0.25*x^4 - x^2 + 1.5};
  \addplot[omThm, dashed, thick, domain=-2.25:2.25, samples=2] {1.0};
  \addplot[omProp, dashed, domain=-2.25:2.25, samples=2] {1.7};
  \node[font=\small, omThm, anchor=west] at (axis cs:-2.2,0.86) {$E_m$};
  \node[font=\small, omProp, anchor=west] at (axis cs:-2.2,1.85) {$E_m'$};
  \node[font=\small, omDef, anchor=west] at (axis cs:1.55,2.1) {$E_p(x)$};
  \fill (axis cs:0.765,1.0) circle (1.6pt) node[font=\small, below=1pt] {$x_1$};
  \fill (axis cs:1.848,1.0) circle (1.6pt) node[font=\small, below right=-1pt] {$x_2$};
  \fill[omMeth] (axis cs:1.414,0.5) circle (2.2pt) node[font=\small, below] {stable};
  \fill[omMeth] (axis cs:0,1.5) circle (2.2pt) node[font=\small, above] {unstable};
\end{axis}
\end{tikzpicture}

{\small An energy diagram, read without solving anything: with energy
$E_m$ the body oscillates between $x_1$ and $x_2$, fastest where the
gap is largest; with $E_m'$ it crawls over the hill and escapes.}
\end{omfigure}

\begin{proposition}[Force from the landscape; equilibria]\label{prop:g12:work-and-energy:equilibria}
The conservative force along $x$ is minus the slope of $E_p$:
$F = -\,\dd E_p/\dd x$ --- it points \emph{downhill} on the diagram,
and vanishes where the tangent is horizontal: an equilibrium. A
minimum of $E_p$ is \emph{stable}\index{stable equilibrium} --- the
displaced body is pushed back, a ball in a valley; a maximum is
\emph{unstable}\index{unstable equilibrium} --- it is pushed away, a
ball on a hilltop.
\end{proposition}

\begin{proof}
On a small step $\dd x$ the force works $F\,\dd x = E_p(x) -
E_p(x + \dd x) = -\dd E_p$; divide by $\dd x$. Downhill on either side
of a valley points back in; around a hilltop, out.
\end{proof}

\begin{remark}[Oscillations and escape]\label{rem:g12:work-and-energy:escape}
The small oscillations of \cref{ch:g12:oscillators-and-time} live at
the bottoms of $E_p$ valleys. And far from Earth gravity weakens: the
$E_p(r)$ curve climbs ever more slowly toward a finite ceiling (the
area under the shrinking force-graph is finite). A rocket whose $E_m$
reaches the ceiling never meets a turning point: it
\emph{escapes}\index{escape speed} --- about \qty{11.2}{km/s} from
the ground (\cref{exo:g12:work-and-energy:13}); below that it is
bound. The ceiling itself is computed in the Year 1 volume.
\end{remark}

\begin{example}[Loop-the-loop]\label{ex:g12:work-and-energy:loop}
A marble released from height $H$ coasts into a vertical loop of
radius $R$. Circular dynamics (\cref{ch:g12:newtons-laws}) demands
$v_{\text{top}}^2 \geq gR$ at the top --- the datum borrowed by last
year's roller-coaster problem, now ours. Conservation from rest at $H$
to the top (height $2R$) gives $v_{\text{top}}^2 = 2g(H - 2R)$, so
$H \geq \tfrac52 R$: half a radius above the loop top, whatever the
mass --- a little more in the real, rubbing world.
\end{example}

\begin{omfigure}
\begin{tikzpicture}[font=\small, scale=0.9]
  \draw[thick] (0,3.0) .. controls (0.9,1.1) and (1.6,0) .. (3.1,0) -- (6.9,0);
  \draw[thick] (5.5,1.2) circle (1.2);
  \fill[omDef] (0,3.0) circle (2.2pt);
  \draw[dashed, omExo] (4.0,2.4) -- (7.6,2.4) node[right] {$2R$};
  \draw[dashed, omThm] (-0.4,3.0) -- (7.6,3.0) node[right] {$H_{\min} = \tfrac52 R$};
  \draw[Stealth-Stealth, omThm] (7.35,2.4) -- (7.35,3.0);
  \node[right, omThm] at (7.35,2.7) {$\tfrac{R}{2}$};
  \draw[-Stealth, omMeth, thick] (5.5,2.4) -- (4.5,2.4) node[above right=0pt and 2pt] {$v_{\text{top}}$};
  \fill (5.5,2.4) circle (1.5pt);
\end{tikzpicture}

{\small Loop-the-loop: the release sits half a radius above the top,
so the leftover $E_k$ keeps the marble pressed on ($v_{\text{top}}^2 = gR$).}
\end{omfigure}

\section{Exercises}

\begin{exercise}[$\star$]\label{exo:g12:work-and-energy:1}
A spring of stiffness $k = \qty{200}{N/m}$ is stretched by
\qty{5.0}{cm}. Compute the stored energy; what does it become if the
stretch doubles? How much does the \emph{second} \qty{5.0}{cm} cost?
\end{exercise}

\begin{exercise}[$\star$]\label{exo:g12:work-and-energy:2}
A \qty{0.90}{kg} drone flies a wandering route from a terrace at
$z = \qty{12}{m}$ to a balcony at $z = \qty{30}{m}$. Work of its
weight? Along the straight climb? On the round trip?
\end{exercise}

\begin{exercise}[$\star$]\label{exo:g12:work-and-energy:3}
Conservative, non-conservative, or workless? Justify from paths:
(a) the weight; (b) the spring force; (c) sliding friction; (d) the
normal force on a sliding box; (e) air drag.
\end{exercise}

\begin{exercise}[$\star$]\label{exo:g12:work-and-energy:4}
A pendulum of length \qty{2.0}{m} is released at rest at
$\qty{60}{\degree}$ from the vertical. Find its speed at the lowest
point. Why does the wire's tension never enter the balance?
\end{exercise}

\begin{exercise}[$\star$]\label{exo:g12:work-and-energy:5}
A cyclist rides at \qty{9.0}{m/s} on the flat, delivering
\qty{250}{W}: what total resistive force is she fighting? Climbing at
\qty{5.0}{m/s} with the same power: what force does she overcome?
\end{exercise}

\begin{exercise}[$\star\star$]\label{exo:g12:work-and-energy:6}
A toy launcher's spring ($k = \qty{450}{N/m}$) is compressed by
\qty{8.0}{cm} behind a \qty{20}{g} dart. Find the stored energy, the
launch speed, and the height reached if fired straight up (no air).
\end{exercise}

\begin{exercise}[$\star\star$]\label{exo:g12:work-and-energy:7}
Stretching a rubber band, a student measures $F = 0$, $3.0$, $7.0$,
$12.0$, \qty{18.0}{N} at $x = 0$, $2.0$, $4.0$, $6.0$, \qty{8.0}{cm}.
Estimate the stored energy by trapezoidal areas; compare with
$\tfrac12 kx^2$, taking $k = F/x$ at the end. Why do they differ?
\end{exercise}

\begin{exercise}[$\star\star$]\label{exo:g12:work-and-energy:8}
A \qty{55}{kg} sledder starts at rest, drops \qty{8.0}{m} along a
\qty{25}{m} curved run, arrives at \qty{9.0}{m/s}. Mechanical energy
lost? Average friction force? Why is the run's exact shape irrelevant?
\end{exercise}

\begin{exercise}[$\star\star$]\label{exo:g12:work-and-energy:9}
A marble track has a loop of radius \qty{20}{cm}. Find the minimum
release height (from rest, no friction). Released from \qty{60}{cm}:
speed at the loop top, and ratio $v_{\text{top}}^2/(gR)$? Why must a
real track be launched from higher still?
\end{exercise}

\begin{exercise}[$\star\star$]\label{exo:g12:work-and-energy:10}
A \qty{100}{g} bead slides without friction along an axis, with
$E_p(x) = 8.0\,x^2$ (joules, $x$ in metres) and $E_m = \qty{2.0}{J}$.
Turning points? Maximum speed? What kind of motion is this?
\end{exercise}

\begin{exercise}[$\star\star$]\label{exo:g12:work-and-energy:11}
A \qty{1200}{kg} car at \qty{2.0}{m/s} is stopped by a bumper spring
of stiffness $k = \qty{6.0e5}{N/m}$. Find the maximum compression.
At \qty{4.0}{m/s}: why not four times as much, the energy quadrupling?
\end{exercise}

\begin{exercise}[$\star\star\star$]\label{exo:g12:work-and-energy:12}
A \qty{0.50}{kg} cart rolls on a track whose $E_p(x)$ has a
\qty{6.0}{J} hilltop ($x = 0$) and a \qty{2.0}{J} valley
($x = \qty{1.5}{m}$). (a) Arriving from the left with
$E_m = \qty{5.0}{J}$: does it pass? What instead? (b) With
$E_m = \qty{7.0}{J}$: speeds at hilltop and valley? (c) Equilibria?
\end{exercise}

\begin{exercise}[$\star\star\star$]\label{exo:g12:work-and-energy:13}
Admitting that hauling a mass $m$ from the ground to arbitrarily far
costs $mgR_E$, $R_E = \qty{6.37e6}{m}$ (the finite area under the
weakening-gravity graph), compute the \emph{escape speed} from Earth.
Why the same for probe and pebble? And a launch just below it?
\end{exercise}

\begin{exercise}[$\star\star\star$]\label{exo:g12:work-and-energy:14}
A \qty{250}{g} glider oscillates on a spring ($k = \qty{40}{N/m}$)
with amplitude \qty{6.0}{cm} (\cref{ch:g12:oscillators-and-time}).
Total energy? Maximum speed? Positions where $E_k = E_p$?
\end{exercise}

\begin{exercise}[$\star\star\star$]\label{exo:g12:work-and-energy:15}
A pole-vaulter sprints at \qty{9.5}{m/s}. Estimate the height his
centre of mass gains if the pole converts all his kinetic energy, and
the bar cleared (centre of mass starting \qty{1.0}{m} up). The record
is about \qty{6.3}{m}: what does the estimate catch, and leave out?
\end{exercise}

\section{Problem: The ski jump}

\begin{problem}[{Weekend problem --- the ski jump: an in-run audited
by a speed gun, a free-fall flight, the gentle geometry of a steep
landing, and the tower a designer must buy once friction bills}]
\label{pb:g12:work-and-energy:1}
An Olympic ski jumper ($m = \qty{70}{kg}$, skis included) starts at
rest from a gate $H = \qty{45}{m}$ above the takeoff table, slides
$L = \qty{90}{m}$ of curved in-run, and leaves the table horizontally;
a speed gun there reads $v_0 = \qty{26}{m/s}$. The landing point lies
$h = \qty{40}{m}$ lower, on a $\qty{35}{\degree}$ slope. Parts II
and III neglect air resistance.

\medskip\noindent\textbf{Part I --- The in-run.}
\begin{enumerate}
  \item Predict the frictionless takeoff speed; why is the curve of
        the in-run irrelevant?
  \item Reference at the table: compute $E_m$ at the gate and at takeoff.
  \item Compute $\Delta E_m$; which theorem names the culprits, and who?
  \item Deduce the average resistive force (snow plus air) over the run.
  \item What fraction of $E_m$ is lost? Why wax the skis and crouch?
\end{enumerate}

\medskip\noindent\textbf{Part II --- The flight.}
\begin{enumerate}[resume]
  \item Only the weight acts in flight: which motion of
        \cref{ch:g12:projectile-motion} is this? Velocity components
        at landing?
  \item Landing speed, twice: from components, and from $E_m$
        conservation. In \unit{km/h}?
  \item Compute the flight time and the horizontal distance flown.
  \item Real jumpers ride the air like a wing, flying much farther:
        can air \emph{raise} the landing speed? Argue with the
        mechanical-energy theorem.
\end{enumerate}

\medskip\noindent\textbf{Part III --- The landing.}
\begin{enumerate}[resume]
  \item Compute the angle of the landing velocity below the horizontal.
  \item Decompose it into components parallel and perpendicular to
        the slope.
  \item The legs absorb only the perpendicular part: compute that
        kinetic energy; compare it with the total.
  \item Flat ground would absorb the lot: to what vertical drop is
        each landing equivalent ($h_{\text{eq}} = v_\perp^2/2g$,
        resp.\ $v^2/2g$)?
  \item One sentence: why are landing hills steep and curved?
\end{enumerate}

\medskip\noindent\textbf{Part IV --- The designer's tower.}
A smaller hill needs takeoff at only $v_0' = \qty{24}{m/s}$; its
straight in-run descends at $\qty{35}{\degree}$, so gate height $H'$
means track length $H'/\sin\qty{35}{\degree}$; same \qty{80}{N} drag.
\begin{enumerate}[resume]
  \item Compute the gate height a frictionless designer would build.
  \item Write the mechanical-energy theorem, gate to table, for $H'$.
  \item Solve for $H'$.
  \item By what percentage did friction raise the tower?
  \item The mass no longer cancels: recompute $H'$ for a
        \qty{60}{kg} jumper. Who needs the taller tower, and why?
  \item Punchline: the tower announced, and friction's share of it.
\end{enumerate}
\end{problem}
